\paragrafos{Σχετικές θέσεις ευθειών στο επίπεδο}{Β.3.7, «Σχετικές θέσεις ευθειών στο επίπεδο»}

\begin{anagke}
Ο παππούς Γιάννης δείχνει πώς θα φυτευτούν οι σειρές με τα καρότα στον κήπο: «Κάθε σειρά θα είναι ευθεία, όλες με την ίδια κατεύθυνση. Καμία δεν θα ακουμπά την άλλη, όσο κι αν τις συνεχίσεις». Η Αγγελική σκέφτεται τις γραμμές στην εθνική οδό και τις ράγες του τρένου: «Δεν συναντιούνται ποτέ, ούτε αν προχωρήσουμε χιλιόμετρα». Ο Ηλίας δείχνει δύο μονοπάτια του κήπου που διασταυρώνονται: «Αυτά ακουμπούν σε ένα σημείο».

«Μα αν συνεχίσεις ένα μονοπάτι που φαινομενικά δεν διασταυρώνεται, μπορεί να διασταυρωθεί αργότερα!», διαμαρτύρεται ο Στέλιος. Η Πολυξένη προτείνει ένα τεστ: «Ας φανταστούμε ότι οι ευθείες δεν τελειώνουν ποτέ, ούτε στο χαρτί. Τότε δύο ευθείες ή έχουν ένα κοινό σημείο ή δεν έχουν κανένα». Η κυρία Αλεξάνδρα προσθέτει μια ερώτηση: «Και πόσες παράλληλες στην ευθεία σας περνούν από ένα σημείο που βρίσκεται δίπλα της;».

Σε αυτή την παράγραφο ονομάζουμε τις δύο δυνατές σχετικές θέσεις δύο ευθειών του επιπέδου, τις τεμνόμενες και τις παράλληλες. Θα μάθουμε και το αίτημα του Ευκλείδη για τις παράλληλες, καθώς και πώς φέρνουμε παράλληλη από ένα σημείο.
\end{anagke}

\begin{lexi}{παράλληλες, τεμνόμενες}{παράλληλες, τεμνόμενες}
Οι \lexiword{παράλληλες} (παρ' άλληλες) πηγαίνουν η μία «δίπλα στην άλλη» χωρίς να συναντιούνται ποτέ. Οι \lexiword{τεμνόμενες} είναι εκείνες που «τέμνονται», δηλαδή «σταυρώνουν». Το κοινό σημείο τους λέγεται \lexiword{σημείο τομής}.
\end{lexi}

\begin{metafrasi}
Δύο ευθείες χωρίς κανένα κοινό σημείο. & παράλληλες: $\varepsilon_1\parallel\varepsilon_2$\\
Δύο ευθείες με ένα κοινό σημείο. & τεμνόμενες\\
Το κοινό τους σημείο. & σημείο τομής $\mathrm{Α}$\\
Τμήματα πάνω σε παράλληλες ευθείες. & $\mathrm{ΑΒ}\parallel\Gamma\Delta$\\
Δύο ευθείες που τέμνονται σε ορθές γωνίες. & κάθετες: $\varepsilon_1\perp\varepsilon_2$\\
\end{metafrasi}

\tmima{Η θεωρία}

\begin{orismos}{παράλληλες και τεμνόμενες ευθείες}
Δύο ευθείες του επιπέδου λέγονται \textbf{παράλληλες} όταν δεν έχουν κανένα κοινό σημείο, όσο κι αν προεκταθούν. Γράφουμε $\varepsilon_1\parallel\varepsilon_2$.\\
Δύο ευθείες του επιπέδου λέγονται \textbf{τεμνόμενες} όταν έχουν ένα μόνο κοινό σημείο. Το σημείο αυτό λέγεται \textbf{σημείο τομής} τους.
\end{orismos}

\begin{center}
\begin{tikzpicture}[scale=0.9]
\draw[thick,{Stealth}-{Stealth}] (0,0)--(4,0); \node[right] at (4,0) {$\varepsilon_2$}; \draw[thick,{Stealth}-{Stealth}] (0,1.3)--(4,1.3); \node[right] at (4,1.3) {$\varepsilon_1$};
\node at (2,-0.6) {$\varepsilon_1\parallel\varepsilon_2$};
\begin{scope}[xshift=7cm]
\draw[thick,{Stealth}-{Stealth}] (0,0)--(4,1.6); \draw[thick,{Stealth}-{Stealth}] (0,1.5)--(4,-0.1); \fill (2,0.77) circle (2pt) node[below] {$\mathrm{Α}$};
\node at (2,-0.6) {τεμνόμενες};
\end{scope}
\end{tikzpicture}
\end{center}

Δύο διαφορετικές ευθείες του ίδιου επιπέδου έχουν δύο μόνο δυνατότητες: ή είναι \emph{τεμνόμενες} ή είναι \emph{παράλληλες}. (Αν είχαν δύο κοινά σημεία, θα ταυτίζονταν, γιατί από δύο σημεία περνά μία μόνο ευθεία.) Δύο τμήματα λέγονται παράλληλα όταν βρίσκονται σε παράλληλες ευθείες.

Οι κάθετες ευθείες είναι μια ειδική περίπτωση τεμνόμενων ευθειών: τέμνονται και σχηματίζουν ορθές γωνίες.

\begin{kanonas}{δύο κάθετες στην ίδια ευθεία}
Δύο ευθείες που είναι κάθετες στην ίδια ευθεία είναι μεταξύ τους \textbf{παράλληλες}.
\end{kanonas}

\begin{center}
\begin{tikzpicture}[scale=0.9]
\draw[thick] (0,0.2)--(5,0.2) node[right] {$\varepsilon$};
\draw[thick] (1,-0.8)--(1,3) node[above] {$\varepsilon_1$}; \draw[thick] (3.6,-0.8)--(3.6,3) node[above] {$\varepsilon_2$};
\draw (1,0.5)--(1.3,0.5)--(1.3,0.2); \draw (3.6,0.5)--(3.9,0.5)--(3.9,0.2);
\node at (2.3,-1.2) {$\varepsilon_1\perp\varepsilon$, $\varepsilon_2\perp\varepsilon$, άρα $\varepsilon_1\parallel\varepsilon_2$};
\end{tikzpicture}
\end{center}

Έτσι φέρνουμε παράλληλη προς μια ευθεία $\zeta$ από ένα σημείο $\mathrm{Α}$ εκτός αυτής: φέρνουμε από το $\mathrm{Α}$ την κάθετη $\mathrm{ΑΒ}$ προς τη $\zeta$ και μετά, στο $\mathrm{Α}$, την κάθετη προς την $\mathrm{ΑΒ}$. Στην πράξη το κάνουμε με γνώμονα και χάρακα: ακουμπάμε τον χάρακα στη δοσμένη ευθεία, τον γνώμονα πάνω στον χάρακα και τον σύρουμε μέχρι να περάσει από το σημείο.

\begin{idiotita}{αίτημα του Ευκλείδη}
Από ένα σημείο εκτός ευθείας διέρχεται \textbf{μία και μόνο μία} ευθεία παράλληλη προς αυτή.
\end{idiotita}

Το αίτημα αυτό το δεχόμαστε χωρίς απόδειξη. Ανήκει στα «Στοιχεία» του Ευκλείδη, που γράφτηκαν στην Αλεξάνδρεια γύρω στο $300$ π.Χ. Από αυτό προκύπτει ότι αν δύο ευθείες είναι παράλληλες προς μια τρίτη, τότε είναι παράλληλες μεταξύ τους: αλλιώς θα είχαν σημείο τομής και από αυτό θα περνούσαν δύο παράλληλες προς την τρίτη.

\tmimasmall{Απόσταση δύο παράλληλων ευθειών}

Αν από ένα σημείο της μίας από δύο παράλληλες ευθείες φέρουμε το κάθετο τμήμα προς την άλλη, το μήκος του είναι η \textbf{απόσταση των παραλλήλων}. Όπου κι αν πάρουμε το σημείο, το κάθετο τμήμα έχει το ίδιο μήκος: οι παράλληλες είναι «ισαπέχουσες».

\begin{center}
\begin{tikzpicture}[scale=0.9]
\draw[thick] (0,0)--(5,0) node[right] {$\varepsilon_2$}; \draw[thick] (0,1.5)--(5,1.5) node[right] {$\varepsilon_1$};
\draw[gblue,thick] (1,0)--(1,1.5) node[midway,left] {\small $\delta$}; \draw[gblue,thick] (3.2,0)--(3.2,1.5) node[midway,right] {\small $\delta$};
\draw (1,0.25)--(1.25,0.25)--(1.25,0); \draw (3.2,0.25)--(3.45,0.25)--(3.45,0);
\end{tikzpicture}
\end{center}

\begin{prosoxi}
\begin{itemize}[leftmargin=1.2em,itemsep=0.3ex,topsep=2pt]
\item Για να είναι παράλληλες δύο ευθείες δεν αρκεί να μην συναντιούνται στο σχέδιο: πρέπει να μην έχουν κοινό σημείο όσο κι αν προεκταθούν. Δύο τμήματα που δεν ακουμπούν μπορεί να βρίσκονται σε τεμνόμενες ευθείες.
\item Δύο παράλληλα τμήματα δεν χρειάζεται να έχουν το ίδιο μήκος.
\item Οι κάθετες είναι τεμνόμενες, αλλά οι τεμνόμενες δεν είναι πάντα κάθετες.
\end{itemize}
\end{prosoxi}

\tmima{Λυμένα παραδείγματα}

\paradeigma{Παράλληλες ή τεμνόμενες;}
\tixero Οι ευθείες: οι ράγες του τρένου, δύο μονοπάτια που διασταυρώνονται, οι γραμμές του τετραδίου.\par
\tithelo Να πω ποιες ευθείες είναι παράλληλες και ποιες τεμνόμενες.\par
\possindeo Παράλληλες είναι όσες δεν έχουν κοινό σημείο, τεμνόμενες όσες έχουν ένα.\par
\apantisi Ράγες και γραμμές του τετραδίου: παράλληλες. Μονοπάτια: τεμνόμενες.\par
\elegxos Οι ράγες δεν θα συναντηθούν ούτε αν προεκταθούν.

\paradeigma{Φέρνω παράλληλη}
\tixero Μια ευθεία $\zeta$ και ένα σημείο $\mathrm{Α}$ εκτός της.\par
\tithelo Η ευθεία $\varepsilon$ από το $\mathrm{Α}$ παράλληλη προς τη $\zeta$.\par
\possindeo Φέρνω από το $\mathrm{Α}$ το κάθετο τμήμα $\mathrm{ΑΒ}$ προς τη $\zeta$. Στο $\mathrm{Α}$ φέρνω την κάθετη στην $\mathrm{ΑΒ}$, δηλαδή την $\varepsilon$. Οι $\zeta$ και $\varepsilon$ είναι κάθετες στην ίδια ευθεία.\par
\apantisi Η $\varepsilon$ είναι παράλληλη προς τη $\zeta$ και είναι η μοναδική που περνά από το $\mathrm{Α}$.\par
\elegxos Η απόσταση των $\varepsilon$ και $\zeta$ είναι $\mathrm{ΑΒ}$, και μετρώντας σε δεύτερο σημείο βρίσκουμε το ίδιο μήκος.

\paradeigma{Δύο παράλληλες σε απόσταση}
\tixero Μια ευθεία $\varepsilon$.\par
\tithelo Ευθείες παράλληλες προς την $\varepsilon$ που απέχουν $3$ cm από αυτήν.\par
\possindeo Παίρνω ένα σημείο $\mathrm{Α}$ της $\varepsilon$ και φέρνω το κάθετο τμήμα $\mathrm{ΑΒ}=3$ cm προς τη μία μεριά και $\mathrm{ΑΒ}'=3$ cm προς την άλλη. Από τα $\mathrm{Β}$, $\mathrm{Β}'$ φέρνω παράλληλες προς την $\varepsilon$.\par
\apantisi Υπάρχουν δύο τέτοιες ευθείες, μία σε κάθε μεριά της $\varepsilon$.\par
\elegxos Και οι δύο έχουν απόσταση $3$ cm από την $\varepsilon$.

\paradeigma{Ευθεία ή τμήμα;}
\tixero Δύο τμήματα $\mathrm{ΑΒ}$ και $\Gamma\Delta$ που δεν ακουμπούν στο χαρτί.\par
\tithelo Είναι παράλληλα;\par
\possindeo Προεκτείνω και τα δύο: αν οι ευθείες τους τέμνονται κάπου εκτός χαρτιού, τα τμήματα δεν είναι παράλληλα.\par
\apantisi Μπορεί να μην είναι: αποφασίζω από τις ευθείες, όχι από τα τμήματα.\par
\elegxos Αν οι προεκτάσεις δεν συναντιούνται, είναι παράλληλα.

\tmima{Ασκήσεις}

\zoni{Το κατάλαβα}{παράλληλες και τεμνόμενες}

\ex Σωστό ή λάθος; α) Δύο ευθείες του επιπέδου που δεν έχουν κανένα κοινό σημείο είναι παράλληλες. \quad β) Δύο τμήματα που δεν έχουν κανένα κοινό σημείο είναι πάντα παράλληλα. \quad γ) Αν δύο ευθείες είναι κάθετες μεταξύ τους, τότε είναι τεμνόμενες. \quad δ) Από ένα σημείο εκτός ευθείας περνά μία μόνο παράλληλη προς αυτήν. \quad ε) Δύο ευθείες κάθετες στην ίδια ευθεία είναι παράλληλες.

\ex Πες αν οι ευθείες είναι παράλληλες ή τεμνόμενες: α) οι δύο γραμμές διαχωρισμού μιας λωρίδας του δρόμου, β) οι ράγες του τρένου, γ) δύο δρόμοι που διασταυρώνονται, δ) οι απέναντι πλευρές ενός θρανίου, ε) τα δύο χέρια του ρολογιού στις $3{:}00$.

\ex Σχεδίασε μια ευθεία $\varepsilon$ και δύο σημεία $\mathrm{Α}$, $\mathrm{Β}$ εκτός της. Από τα $\mathrm{Α}$ και $\mathrm{Β}$ φέρε παράλληλες προς την $\varepsilon$. Τι παρατηρείς; Πότε ταυτίζονται οι δύο ευθείες;

\ex Σχεδίασε μια ευθεία $\varepsilon$ και ένα σημείο $\mathrm{Α}$ εκτός της. Φέρε από το $\mathrm{Α}$ την κάθετη προς την $\varepsilon$ και μετά, στο $\mathrm{Α}$, την κάθετη προς αυτήν. Τι σχέση έχει η δεύτερη κάθετη με την $\varepsilon$;

\ex Σχεδίασε μια γωνία $x\hat{\mathrm{Ο}}y=50^\circ$. Από ένα σημείο $\mathrm{Α}$ της $\mathrm{Ο}x$ φέρε παράλληλη προς την $\mathrm{Ο}y$ και από ένα σημείο $\mathrm{Β}$ της $\mathrm{Ο}y$ παράλληλη προς την $\mathrm{Ο}x$. Αν $\mathrm{Κ}$ είναι το σημείο τομής τους, μέτρησε τη γωνία $\mathrm{ΑΚΒ}$.

\ex Σχεδίασε μια ευθεία $\varepsilon$ και δύο ευθείες που απέχουν $4$ cm από αυτήν και είναι παράλληλες προς αυτήν. Πόσες τέτοιες ευθείες υπάρχουν;

\ex Σχεδίασε δύο τμήματα που δεν τέμνονται πάνω στο χαρτί, αλλά οι ευθείες τους τέμνονται εκτός χαρτιού. Είναι παράλληλα τα τμήματα;

\zoni{Το χρησιμοποιώ}{κήποι, ράγες, πλακάκια}

\ex \textbf{Οι σειρές των καρότων.} Ο παππούς Γιάννης φυτεύει πέντε παράλληλες σειρές καρότα και κάθε δύο διαδοχικές απέχουν $40$ cm. Πόση είναι η απόσταση ανάμεσα στην πρώτη και στην τελευταία σειρά;

\ex \textbf{Οι ράγες.} Οι ράγες του τρένου απέχουν $1{,}5$ μ. σε ένα σημείο. Πόσο απέχουν $2$ χιλιόμετρα πιο κάτω; Εξήγησε.

\ex \textbf{Οι δρόμοι.} Οι οδοί $\varepsilon_1$ και $\varepsilon_2$ είναι κάθετες στην οδό $\varepsilon_3$. Πώς είναι οι $\varepsilon_1$ και $\varepsilon_2$ μεταξύ τους;

\ex \textbf{Τα πλακάκια.} Η θεία Μαρία στρώνει πλακάκια στην ταράτσα. Οι αρμοί προς τη μία κατεύθυνση είναι κάθετοι στους αρμούς της άλλης κατεύθυνσης. Πώς είναι οι αρμοί της ίδιας κατεύθυνσης μεταξύ τους; Εξήγησε.

\ex \textbf{Το τετράδιο.} Πώς θα σχεδιάσεις με χάρακα και γνώμονα μια ευθεία παράλληλη προς μια δοσμένη, που να περνά από ένα δοσμένο σημείο; Γράψε τα βήματα.

\zoni{Το σκέφτομαι}{θέσεις και αιτήματα}

\ex Γιατί δύο διαφορετικές ευθείες του επιπέδου δεν μπορούν να έχουν δύο κοινά σημεία;

\ex Αν $\varepsilon_1\parallel\varepsilon_2$ και $\varepsilon_2\parallel\varepsilon_3$, τι σχέση έχουν οι $\varepsilon_1$ και $\varepsilon_3$; Εξήγησε με το αίτημα του Ευκλείδη.

\ex Από ένα σημείο εκτός ευθείας $\varepsilon$ πόσες ευθείες περνούν που δεν τέμνουν την $\varepsilon$; Και πόσες που την τέμνουν;

\ex Είναι κάθετες οι τεμνόμενες ευθείες; Είναι τεμνόμενες οι κάθετες ευθείες; Δώσε παραδείγματα.

\zoni{Για το σπίτι}{ψάξε παράλληλες}

\ex Βρες στο σπίτι σου πέντε ζεύγη παράλληλων ευθειών (ράφια, πλευρές παραθύρου κ.λπ.) και δύο ζεύγη κάθετων.

\ex Δίπλωσε ένα φύλλο στη μέση. Δίπλωσέ το ξανά ώστε η πτυχή να πέσει πάνω στον εαυτό της. Έτσι έχεις φτιάξει κάθετες πτυχές. Φτιάξε με τον ίδιο τρόπο δεύτερη κάθετη στην ίδια πτυχή. Πώς είναι οι δύο μεταξύ τους;
